% GENERATED FILE. Edit canonical lessons, then run npm run generate:books.
% canonical-source-hash: fnv1a64:859e193b658bdd3d
% canonical-lessons: MR-C342-hetu, MR-C342-uddes, MR-C342-parinam, MR-C342-spastikaran, MR-C342-nirman-hone

\chapter{Why Things Happen}
\label{ch:mr-why-things-happen}
\hlchaptermodality{342}

\begin{chapteropening}
\textbf{By the end of this chapter:} \emph{I can talk about motives, purposes and results.}

Complete the last lesson of chapter 342: I can talk about motives, purposes and results.
\end{chapteropening}

\section[\texorpdfstring{\mr{हेतू} (hetū)}{हेतू (hetū)}]{\mr{हेतू} (hetū) — a motive; an intention}

\label{lesson:MR-C342-hetu}



\begin{quote}
Before the new one: say the Marathi for dance, then the Marathi for a stage.
\end{quote}

\subsection*{You'll want to know: \mr{हेतू}}
\textbf{\mr{हेतू}} — \emph{hetū} — \textquotedblleft{}a motive; an intention\textquotedblright{}.

\begin{quote}
\mr{त्याच्या} \mr{मदतीमागचा} \mr{हेतू} \mr{चांगला} \mr{होता}.
\end{quote}

\begin{quote}
\emph{tyācyā madtīmāgcā hetū cāṅglā hotā.}
\end{quote}

\begin{quote}
The motive behind his help was good.
\end{quote}

Masculine.

\mr{हेतूने}, with the intention: \mr{शिकण्याच्या} \mr{हेतूने}, with the aim of learning.

\begin{cousinweb}
\textbf{\mr{हेतू}} (\emph{hetū}) is a tatsama, taken from Sanskrit \emph{hetu}, \textquotedblleft{}impulse, motive, cause\textquotedblright{} (Monier-Williams). Marathi lengthens the final vowel.
\end{cousinweb}

\begin{grammarlens}[title={What you've built}]
One more word for giving reasons and opinions.
\end{grammarlens}

\subsection*{Your turn}
\begin{itemize}
\raggedright
  \item \emph{Say it:} \emph{hetū}
  \item \emph{Say it:} \emph{hetū}, once more
  \item \emph{Say it:} \emph{tyācyā madtīmāgcā hetū cāṅglā hotā.} — The motive behind his help was good.
  \item \emph{Recall:} say the Marathi for dance, then the Marathi for a stage, then say \emph{hetū} again
\end{itemize}

\begin{tcolorbox}[breakable,colback=teal!4,colframe=teal!35!black,title={Before you move on}]
What does \mr{हेतू} mean? (\textquotedblleft{}A motive; an intention\textquotedblright{}.) Say it once more.
\end{tcolorbox}
\section[\texorpdfstring{\mr{उद्देश} (uddeś)}{उद्देश (uddeś)}]{\mr{उद्देश} (uddeś) — a purpose; an aim}

\label{lesson:MR-C342-uddes}



\begin{quote}
Before the new one: say the Marathi for a stage, then the Marathi for a motive.
\end{quote}

\subsection*{You'll want to know: \mr{उद्देश}}
\textbf{\mr{उद्देश}} — \emph{uddeś} — \textquotedblleft{}a purpose; an aim\textquotedblright{}.

\begin{quote}
\mr{या} \mr{प्रवासाचा} \mr{उद्देश} \mr{काय} \mr{आहे}?
\end{quote}

\begin{quote}
\emph{yā pravāsācā uddeś kāy āhe?}
\end{quote}

\begin{quote}
What is the purpose of this trip?
\end{quote}

Masculine.

The thing whose purpose it is takes -\mr{चा}: \mr{प्रवासाचा} \mr{उद्देश}.

\begin{cousinweb}
\textbf{\mr{उद्देश}} (\emph{uddeś}) is a tatsama, taken from Sanskrit \emph{uddeśa}, \textquotedblleft{}pointing at; aim, object\textquotedblright{} (Monier-Williams).
\end{cousinweb}

\begin{grammarlens}[title={What you've built}]
One more word for giving reasons and opinions.
\end{grammarlens}

\subsection*{Your turn}
\begin{itemize}
\raggedright
  \item \emph{Say it:} \emph{uddeś}
  \item \emph{Say it:} \emph{uddeś}, once more
  \item \emph{Say it:} \emph{yā pravāsācā uddeś kāy āhe?} — What is the purpose of this trip?
  \item \emph{Recall:} say the Marathi for a stage, then the Marathi for a motive, then say \emph{uddeś} again
\end{itemize}

\begin{tcolorbox}[breakable,colback=teal!4,colframe=teal!35!black,title={Before you move on}]
What does \mr{उद्देश} mean? (\textquotedblleft{}A purpose; an aim\textquotedblright{}.) Say it once more.
\end{tcolorbox}
\section[\texorpdfstring{\mr{परिणाम} (pariṇām)}{परिणाम (pariṇām)}]{\mr{परिणाम} (pariṇām) — a consequence; an effect}

\label{lesson:MR-C342-parinam}



\begin{quote}
Before the new one: say the Marathi for a motive, then the Marathi for a purpose.
\end{quote}

\subsection*{You'll want to know: \mr{परिणाम}}
\textbf{\mr{परिणाम}} — \emph{pariṇām} — \textquotedblleft{}a consequence; an effect\textquotedblright{}.

\begin{quote}
\mr{पावसाचा} \mr{परिणाम} \mr{भाजीच्या} \mr{दरावर} \mr{झाला}.
\end{quote}

\begin{quote}
\emph{pāvsācā pariṇām bhājīcyā darāvar jhālā.}
\end{quote}

\begin{quote}
The rain had an effect on vegetable prices.
\end{quote}

Masculine.

X \mr{चा} \mr{परिणाम} Y \mr{वर} \mr{होतो}: X has an effect on Y.

\begin{cousinweb}
\textbf{\mr{परिणाम}} (\emph{pariṇām}) is a tatsama, taken from Sanskrit \emph{pariṇāma}, \textquotedblleft{}change, development, result\textquotedblright{} (Monier-Williams).
\end{cousinweb}

\begin{grammarlens}[title={What you've built}]
One more word for giving reasons and opinions.
\end{grammarlens}

\subsection*{Your turn}
\begin{itemize}
\raggedright
  \item \emph{Say it:} \emph{pariṇām}
  \item \emph{Say it:} \emph{pariṇām}, once more
  \item \emph{Say it:} \emph{pāvsācā pariṇām bhājīcyā darāvar jhālā.} — The rain had an effect on vegetable prices.
  \item \emph{Recall:} say the Marathi for a motive, then the Marathi for a purpose, then say \emph{pariṇām} again
\end{itemize}

\begin{tcolorbox}[breakable,colback=teal!4,colframe=teal!35!black,title={Before you move on}]
What does \mr{परिणाम} mean? (\textquotedblleft{}A consequence; an effect\textquotedblright{}.) Say it once more.
\end{tcolorbox}
\section[\texorpdfstring{\mr{स्पष्टीकरण} (spaṣṭīkaraṇ)}{स्पष्टीकरण (spaṣṭīkaraṇ)}]{\mr{स्पष्टीकरण} (spaṣṭīkaraṇ) — an explanation}

\label{lesson:MR-C342-spastikaran}



\begin{quote}
Before the new one: say the Marathi for a purpose, then the Marathi for a consequence.
\end{quote}

\subsection*{You'll want to know: \mr{स्पष्टीकरण}}
\textbf{\mr{स्पष्टीकरण}} — \emph{spaṣṭīkaraṇ} — \textquotedblleft{}an explanation\textquotedblright{}.

\begin{quote}
\mr{उशिराबद्दल} \mr{त्याने} \mr{स्पष्टीकरण} \mr{दिलं}.
\end{quote}

\begin{quote}
\emph{uśirābaddal tyāne spaṣṭīkaraṇ dilaṁ.}
\end{quote}

\begin{quote}
He gave an explanation for being late.
\end{quote}

Neuter.

\mr{स्पष्टीकरण} \mr{देणे}, to give an explanation; it is built on \mr{स्पष्ट}, clear.

\begin{cousinweb}
\textbf{\mr{स्पष्टीकरण}} (\emph{spaṣṭīkaraṇ}) is a tatsama, taken from Sanskrit \emph{spaṣṭīkaraṇa}, \textquotedblleft{}making clear\textquotedblright{} (Monier-Williams).
\end{cousinweb}

\begin{grammarlens}[title={What you've built}]
One more word for giving reasons and opinions.
\end{grammarlens}

\subsection*{Your turn}
\begin{itemize}
\raggedright
  \item \emph{Say it:} \emph{spaṣṭīkaraṇ}
  \item \emph{Say it:} \emph{spaṣṭīkaraṇ}, once more
  \item \emph{Say it:} \emph{uśirābaddal tyāne spaṣṭīkaraṇ dilaṁ.} — He gave an explanation for being late.
  \item \emph{Recall:} say the Marathi for a purpose, then the Marathi for a consequence, then say \emph{spaṣṭīkaraṇ} again
\end{itemize}

\begin{tcolorbox}[breakable,colback=teal!4,colframe=teal!35!black,title={Before you move on}]
What does \mr{स्पष्टीकरण} mean? (\textquotedblleft{}An explanation\textquotedblright{}.) Say it once more.
\end{tcolorbox}
\section[\texorpdfstring{\mr{निर्माण} \mr{होणे} (nirmāṇ hoṇe)}{निर्माण होणे (nirmāṇ hoṇe)}]{\mr{निर्माण} \mr{होणे} (nirmāṇ hoṇe) — to come about; to arise}

\label{lesson:MR-C342-nirman-hone}



\begin{quote}
Before the new one: say the Marathi for a consequence, then the Marathi for an explanation.
\end{quote}

\subsection*{You'll want to know: \mr{निर्माण} \mr{होणे}}
\textbf{\mr{निर्माण} \mr{होणे}} — \emph{nirmāṇ hoṇe} — \textquotedblleft{}to come about; to arise\textquotedblright{}.

\begin{quote}
\mr{या} \mr{निर्णयामुळे} \mr{नवा} \mr{प्रश्न} \mr{निर्माण} \mr{झाला}.
\end{quote}

\begin{quote}
\emph{yā nirṇayāmuḷe navā praśna nirmāṇ jhālā.}
\end{quote}

\begin{quote}
A new question arose because of this decision.
\end{quote}

Past (he …): \mr{तो} \mr{निर्माण} \mr{झाला}, \emph{to nirmāṇ jhālā}.

Intransitive: what arises is the subject. \mr{निर्माण} \mr{करणे} is to create.

\begin{cousinweb}
\textbf{\mr{निर्माण} \mr{होणे}} (\emph{nirmāṇ hoṇe}) joins the tatsama \mr{निर्माण}, Sanskrit \emph{nirmāṇa}, \textquotedblleft{}measuring, making\textquotedblright{} (CDIAL 7368), to \mr{होणे}, \textquotedblleft{}to become\textquotedblright{}: to come into being.
\end{cousinweb}

\begin{grammarlens}[title={What you've built}]
One more word for giving reasons and opinions.
\end{grammarlens}

\subsection*{Your turn}
\begin{itemize}
\raggedright
  \item \emph{Say it:} \emph{nirmāṇ hoṇe}
  \item \emph{Say it:} \emph{nirmāṇ hoṇe}, once more
  \item \emph{Say it:} \emph{yā nirṇayāmuḷe navā praśna nirmāṇ jhālā.} — A new question arose because of this decision.
  \item \emph{Recall:} say the Marathi for a consequence, then the Marathi for an explanation, then say \emph{nirmāṇ hoṇe} again
\end{itemize}

\begin{tcolorbox}[breakable,colback=teal!4,colframe=teal!35!black,title={Before you move on}]
What does \mr{निर्माण} \mr{होणे} mean? (\textquotedblleft{}To come about; to arise\textquotedblright{}.) Say it once more.
\end{tcolorbox}
